\section{What the algebra decides}
\label{sec:assessment}

We sort the results of this paper into the four classes announced in \cref{sec:deflation}, and then name the choices that remain. The test applied throughout is the
deflationary one of \cref{cor:nopredict}: a statement about the arena carries information only if it could have come out otherwise, and compatible representation
data is not a prediction.

\subsection{Fixed by the algebra}

At one constituent the algebra fixes more than one might expect. The signature rule separates the bivectors that can gauge from those that can only grade, and the
pole operators grade the algebra by $\Z_2^3$ through a maximally split Cartan subalgebra. The pairing of the two triples has exactly two square roots, the grading root
$\cK$ and the binding root $G$, related by the Cartan involution $\Ad_{e_1e_2e_3}$. Under permutation symmetry the clock is the barycentre of the Fisher--Rao octant, and
exactly that clock makes the conversion invisible in the equations of motion. The imaginary unit is selected, not adjoined: inside the joint commutant of the interior
Lorentz action and the internal bivectors the only complex structures commuting with the gamma matrices are $\pm J$. Charge conjugation is the surplus plane, time reversal
squares to $-\id$, the kinetic forms alone fix the massless gauge algebra $\uu1_\Omega\oplus\su2$, the mass form cuts it to $\uu1_J$, and that $\uu1$ is the whole internal
centraliser. Along the tower the internal centraliser is $\so{3n-1,3n-3}$, nested and over one Lorentz frame, and the diagonal Lorentz action gives a compact internal algebra
forming a reductive dual pair with the Lorentz algebra, with the exact identity $C_{\rm cop}+C_\Delta=\tfrac12n(n+2)$.

In the field-strength reading the centraliser of $G$ is a Lorentz algebra under which $\R^{3,3}$ \emph{is} $\Lambda^2\R^{1,3}$, $G$ is the Hodge star, the two Maxwell
invariants are the real and imaginary parts of one holomorphic quadratic form, and covariance fixes the form of Maxwell's equations through a multiplicity count. In the
interior dynamics, factorisation and binding coexist exactly in neutral signature, and the force law closes in one way only, into a rotor acting by $G$ whose free motion is a
pendulum in the duality angle. In the curved arena the curvature operator lives on the frame, the first Bianchi identity is the single trace $\Tr(G\cR)=0$, and the duality
rotation of the curvature ends, where conservation and variationality both select it, at the double dual contracting to the Einstein tensor; in first-order form the same
rotation is a Barbero--Immirzi angle. Kinematically, choosing a Lorentzian four-plane is choosing a complex structure, and at one constituent the light cone is the set of
chiral spinors squared with the help of $J$.

\subsection{Constructed}

The constructions are existence results; in each case the algebra admits the object and nothing selects it among alternatives. At three constituents the axial projector
$P_*$ of rank $64$ carries the Pati--Salam algebra and one chiral family with a right-handed neutrino, and settles the question of colour with chirality. The diagonal Dirac
block carries a gauge-invariant real $\Cl(1,3)$ built from a copy-singlet cubic. At the mixed second rung the univectorial potential, a Dirac block and the frame field can
be coupled into two-flavour spinor electrodynamics. On the $272$ mass kernels of the chiral block the renormalizable potential of the $(\mathbf{10},\mathbf1,\mathbf3)$ component
is classified, the conditions under which the Standard-Model orbit is its global minimum are found, and the radiative spectrum of all $271$ angular modes is computed, with
the electroweak stage, the seesaw and the Georgi--Jarlskog factor as consequences of a bidoublet vacuum value. At the second rung a Lorentz-invariant, uncharged, non-rigid
vacuum with a charged-scalar current exists on a compact $S^4$ of projectors, with a curved extension in which it is parallel. Both readings of $\R^{3,3}$ construct the
Einstein tensor of a given four-metric.

\subsection{Obstructed}

Each obstruction holds under hypotheses stated where it was proved, and those hypotheses are part of the result. Hypercharge cannot be read off anomaly cancellation for
charges built from three commuting involutions; $\su3$ fits in the effective algebra of no axial projector of the five-projector family; the two central $\uu1$s of $P_*$
cannot be gauged on its matter; the diagonal sector is vector-like on its complete blocks. The two readings of the same six directions exclude one another, and the
obvious repair at the second rung costs the whole internal algebra. The positive-form interior oscillates only at zero momentum. No Lie--Poisson closure lives on the frame
alone, and no inertia is both Lorentz invariant and rate-conserving. Non-orthogonality of $\bE$ and $\bB$ breaks $P$ and $T$, is never spontaneous, and generates no mass.
The paired-mode projector sources only a massive vector, and at the first rung a vacuum that is both massless and Lorentz invariant is rigid. The normal $U(1)$ of the curved
model fails as electromagnetism on three independent counts: a spinless particle cannot see it, its adapted generator does not commute with the physical generators, and it
gravitates with the wrong sign. Ricci flow does not transfer to split signature. No vacuum value of a field built from the internal bivectors and $J$ makes the matter
chiral, Hermitian order parameters keep rank five, and vacuum alignment does not see chirality. In the causal analysis, the six-dimensional free field is ill-posed and has
no time orientation, no $SO(3,3)$-invariant scalar potential of degree below eight isolates the Lorentzian type, every such potential is degenerate with a Lorentz-violating
orbit and has six unprotected flat directions, every invariant kinetic term has ghosts, and every vector from the normal plane is a ghost.

\subsection{Open, and the choices that remain}

\Cref{tab:choices} lists the choices that the algebra leaves free, with the place where each enters. Most of the open problems are downstream of them.

\begin{table}[t]
\centering
\small
\begin{tabular}{@{}>{\raggedright\arraybackslash}p{4.0cm}>{\raggedright\arraybackslash}p{10.4cm}@{}}
\toprule
choice & status\\
\midrule
the projector & a field with a computable configuration space; at rank $64$ at least two admissible orbits, of dimensions $34$ and $20$; a spectrum-only potential is constant on
them, so selection needs couplings that distinguish orbits\\
the Lorentz action & spectator and diagonal actions are not conjugate and trade against each other: a large internal algebra with ghost-signed directions and a chiral block,
or a small compact algebra on vector-like matter\\
the hypercharge direction & a compatible embedding exists on $P_*$ and is output by the vacuum of the kernel field, but the vacuum itself is selected only on an open set of
couplings\\
the binding operator and its reading & $G$ is adopted; conditional uniqueness once its Lorentz algebra is taken, not selection; the field-strength reading and the interior reading
of the same six directions exclude one another\\
scales and inertias & no ratio among the rotor inertia, the frame inertia and the coupling scales is fixed; the inertia of the closure has $21$ free directions satisfying neither
requirement\\
the action & no parent action is derived; the gauge field, the metric and the kernel field are given dynamics by adopted actions whose equations follow\\
the four-plane and positivity & the selector is $J$; the type is not selected by invariant dynamics, and every healthy construction used a positive inner product\\
families & the block carries one family; three copies of a module are a multiplicity index, not three constituents, and nothing reduces the copy multiplicity to three\\
\bottomrule
\end{tabular}
\caption{The choices the algebra leaves open.}
\label{tab:choices}
\end{table}

Several problems are sharply posed by the results. A coupling that distinguishes the points of the compact family of second-rung vacua, done at $n\ge3$ where colour
survives, would turn the photon construction into a selection. A finite search over non-diagonal pairings of the two Lorentz bases at the second rung and above would settle
whether the two readings of $\R^{3,3}$ can share one theory at an acceptable price. The spectrum and the index of the gauge-invariant Dirac operator of the diagonal block on
a curved base are ordinary index-theorem calculations, and a nonvanishing index would not contradict vector-likeness. The global structure of the potential on all $272$ mass
kernels, the choice between the Standard-Model vacuum and its mirror, the electroweak hierarchy, and the explicit breaking of the left-handed fermion number that the axion
bound requires are all open. A mechanism placing the interior momenta on the null cone or on the balance surface, the only places where a constant rate and a bounded frame
coexist, and the quantisation of the rotor, whose angle would then have to be related to an external time, are open. Whether the antisymmetric part of the shift is the
contorsion of an independent connection is a conjecture that is not yet type-correct. A gravitational action relating metric, pairing, connection and projector, from which
the metric would follow rather than be given, is the largest open problem of the curved arena; and in the causal analysis a gauge mechanism for type selection and a principle
that makes the normal translations gauge are the two steps from which a derivation of causal order would have to begin.

\subsection{The test, applied}

The largest gap is the same in every sector, and it is one of selection rather than absence. The gauge algebra is derived and the gauge field is not; the arena is derived
and its dynamics is not; the Einstein tensor of a given metric is constructed and the metric is not determined; the projector is a field whose equation follows once an action
is adopted, but no action is derived and none of its admissible orbits is singled out; the vacuum of the kernel field is selected, but only on a stated open set of couplings
and with field content given. Recovering the hypercharge table, the anomaly conditions or $\sin^2\theta_W=\tfrac38$ is a consistency result, not evidence, since any construction
with that charge content gives them. New mathematics would be a theorem that removes one of the choices of \cref{tab:choices} without adding another; new physics would require
a selected spectrum, a coupling or scale relation, a mass restriction, or a deviation that survives re-expression in the complex language of \cref{thm:realification}. None of
these is claimed here. What the paper does claim is a precise map of where the split-signature Clifford algebra decides and where it leaves the decision to the physicist,
and on that map the decided territory is larger, and its borders sharper, than one might have guessed from the size of the arena.
